\chapter{中世纪手稿选读}
本章节主要介绍一些基础的中世纪手稿判读知识。在之前的章节中，我们采用的都是经过标准正字法修正的文稿。编者在组织这些文字时，有意将它们向现代的冰岛语正字法靠拢以方便阅读。现在，我们将直接阅读一部分原始手稿，手稿中的记载和我们之前看到的“标准化”的文本有几个主要的不同：
\begin{enumerate}
    \item 字体不同。我们将认识加洛林手写体（Carolingian minuscule）及其变体。

    \item 大量省略、缩写、连写。为了节省纸张，这是中世纪手稿的特点。我们将认识一些常见的速写记号。

    \item 粗糙的正字法。在中世纪手稿中，许多发音上的细节并没有体现到文字中，例如字母不分清浊长短等。
\end{enumerate}

\section{基础知识}
\subsection{字体}
绝大多数古诺尔斯语的手稿的成书时间都在1100年后，当时在欧洲大陆流行的字体是所谓的“加洛林体”，这是一种清晰、易于掌握的古拉丁文小写体，由统治法兰克王国的加洛林王朝最先推广，也因此得名。

加洛林体主要是用于书写拉丁文的，但伴随着宗教和文法的影响，日耳曼地区也开始使用这种字体。当然，随之而来的问题就是拉丁字母并不能表示日耳曼语中所有的音，例如擦音\th 和\dh 就没有对应的字母。英国人最早解决了这个问题，他们重新结合了卢恩字母和拉丁字母补充了一部分拉丁语没有的字母，并且改进出了一套适合他们语言的字体：岛屿体（Insular script）。古英语的音系和古诺尔斯语有许多相似之处，于是，当时的北欧人从英国人那里借来了表示日耳曼语特有音素的字母，并对加洛林体做了改动，这种字体后来称为哥特体的前身，称为`praegothica'，即“早期哥特体”。早期哥特体相比加洛林体而言横向间距更加紧致，字母仅靠在一起，连写和合写字母很多，笔锋尖锐。

下图从上到下分别是加洛林体、岛屿体和早期哥特体的样例（理想情况），当然由于中世纪的正字法问题以及各种字体存在的变体，这还不是我们手稿中看到的情况。读者通过此图大致了解各字体的基本样貌即可，请不要把它当成某种标准：
\begin{figure}[htbp]
    \centering
    \includegraphics[width=0.8\textwidth]{figure/script.png}
\end{figure}

在下一节的手稿中，我们遇到的主要是早期哥特体，一些字形和上图中的情况还有一定的区别，值得注意的几点是：
\begin{info}
    \begin{enumerate}[leftmargin=*]
        \item i/j、u/v区分不明显。准确来说此时的手稿还不区分i/j以及u/v, i和j都写作ı，类似于希腊字母iota；u和v都写作v，但底部不如v那么尖锐。有时ı会写到字母右上角；

        \item d写作{\medieval ꝺ}，和岛屿体类似；

        \item f写作{\medieval ꝼ}，它和现代字母的区别在于左边的一竖延申到基准线下方；

        \item s总是写作{\medieval ſ}，这称为`long s'，和l容易搞混；

        \item t写作{\medieval ꞇ}，和c类似，这也是从岛屿体中借来的；

        \item \th 和 \dh 基本和现在的字母形状相似，其中\dh 实际上是在{\medieval ꝺ}的基础上加了一条横杠。不过注意，现代字母中这条横杠是直的，手稿中这条杠在收笔时常向左下方弯曲，形成一个钩子状的符号。这个钩子有时候会靠到{\medieval ꝺ}上，形成一个小圆圈；
    \end{enumerate}
\end{info}
\begin{info}
    \begin{enumerate}[leftmargin=*, start=7]

        \item 从岛屿体中借来了字母{\medieval ꝩ}，这个字母叫作`vend'，由卢恩字母演变而来，音值相当于v或者u，它主要用在词首；

        \item y写作\.y，顶上的一点使之与{\medieval ꝩ}更好地区分开。
    \end{enumerate}
\end{info}
\subsection{合体字}
在中世纪的手稿中，有时会把多于一个字母的合成一个字形，这称为合体字（Ligature）。最典型的例子就是英语字母W，它就是VV或者UU的合体字。有两个合体字比较重要：
% \setmainfont{[Cardo104s.ttf]}
\begin{enumerate}[itemindent=1em]
    \item  {\andron ꜹ}

          这是a和u/v的合体字，如`auka'写作 {\andron ꜹka}.

    \item {\andron {\andron ꝛ}}

          这个字母称为半R（R rotunda）。它一开始是和`o'合写的，如{\andron o{\andron ꝛ}igin}，后来，它也可以和任何右半部分是圆型的字母合写，如`b, p'等。
\end{enumerate}

\subsection{速写记号}
为了加快书写速度和节省纸张，许多单词被缩写、省略或用记号代替，本书把这些记法统称为速写记号（Scribal abbreviation）。一些通用的记号如下表所示，很多记号表达的含义在各种语言中是通用的：
\begin{table}[H]
    \centering
    \begin{tabular}{@{}lll@{}}
        \toprule
        记号                           & 意义       & 说明                                                                                          \\
        \midrule
        {\medieval ħ}                  & hann       & 快速缩写人称代词                                                                              \\
        {\medieval ꝥ}                  & \th at     & 快速缩写指示代词                                                                              \\
        {\medieval ꝁ}                  & konungr    & 快速缩写专有名词                                                                              \\
        {\medieval ⁊}                  & ok         & 形如7，通用的速记符号，表示“和”                                                               \\
        \textarm{m}                    & maðr       & 用卢恩字母的名字快速缩写常见名词                                                              \\
        {\medieval mͬ}                  & maðr       & 用首尾字母快速缩写常见名词                                                                    \\
        {\medieval mꝫ/m;}              & með        & {\medieval ꝫ}是拉丁语中用来缩写et的字母，在古诺尔斯语中表示相近的音值eð，;和{\medieval ꝫ}形近 \\
        {\textvbaraccent{\medieval ꞇ}} & til        & 实际上相当于在t上方简写了一个i                                                                \\
        {\medieval ſ\l}                & skal       & 由ſ向l连笔，省略了中间两个字母                                                                \\
        {\medieval ◌̄}                  & m/n        & 上横杠，表示一个含鼻音的音节，如n\=ea `nema'                                                  \\
        {\medieval ◌̾/◌͛/◌}             & r/ir/er/ri & 闪电状或类似于7的符号，表示一个含r的音节，常见于词尾，如{\medieval marg͛ } `margir'            \\
        {\medieval ◌ᷓ}                  & ra/ar      & 形如w的符号，实际上是加洛林体字母a变体 ({\medieval })，表示ar或/ra，如{\medieval beᷓ} `bera'  \\
        \bottomrule
    \end{tabular}
\end{table}

\subsection{正字法}
了解了字体和记号之后，读者还需要对中世纪的正字法（更准确地来说，是书写习惯）有基本的认识：
\begin{info}
    \begin{enumerate}[leftmargin=*]
        \item 长音。中世纪手稿大多数不标记长音，用`ˊ'标记长音的情况罕见；

        \item 双辅音。常用大写字母表示双辅音，但这个大写字母和小写字母等大。如le\textsc{g}ja = leggja,

        \item 弱读辅音。非重读的/u/和/i/基本写作o和e，因此手稿中的词尾和教材中的略有区别；

        \item 变元音。变元音的记号在手稿中比较混乱，大致有以下几点规律：
              \begin{enumerate}
                  \item ø可写作eo, ø, o, ey等；
                  \item ǫ可写作ǫ, ao, au, o等；
                  \item æ, e可写作e, ę, æ等;
              \end{enumerate}

        \item 常用c代表k音。一般和元音连用时比较常见，如oc = ok, calla = kalla,

        \item \th 和 \dh 常混用。
    \end{enumerate}
\end{info}

\section{例文}
下面的例文摘录自灰天鹅法典（Grágás, Grey Goose Laws），这是13世纪中期的一本法典。手稿节选自冰岛学者Árni Magnússon整理的AM 334 fol系列\footnotemark。
\footnotetext{扫描件可见\href{Handrit.is}{https://myndir.handrit.is/file/Handrit.is/AM\%20334\%20fol./52/SECONDARY\_DISPLAY}}
下面节选的法案（AM 334 fol :: 24r）属于“赡养法”（Ómaga Bálkr\footnotemark）的部分，是关于代养孩子的：
\footnotetext{ómagi或úmagi指的是不能独立生活的人，如小孩、老人、贫民等。法律规定自由民有赡养他们的义务。}
\begin{multicols}{2}
    \begin{figure}[H]
        \includegraphics[width=0.5\textwidth,height=0.73\textwidth]{figure/mss.png}
    \end{figure}
    \newcolumn
    \quad
    \fontsize{14pt}{14pt}\selectfont
    {\medieval
    \hfill \Eclosedunical ꝼ \\
    \noindent mͬ ꝼøþ͛ barn {\andron ꜹ}ð{\andron ꝛ}\=o \=me oc ꞇekr \vbart ꝼoſꞇrſ oc ſ\l\\
    ħ ꝥ barn ꝼøða þar \vbart er ꝥ er \textsc{xvi} ꝩᷓ gam\\
    alꞇ. \Eunical \nwithdescender \ eꝼ barne ꝩ͛þa leſꞇır aþ{\andron ꝛ} ꝥ er \textsc{xvi} ꝩᷓ\\
    gamallꞇ, þa ſ{\l} ꝼeeꞇ apꞇr g\=aga leıgo l{\andron ꜹ}ſꞇ.\\
    oc ꞇelıa ı\=aꞇ \vbart hv͛ſ ꝩeꞇᷓr ı\nwithdescender ſꞇøða\nwithdescender . \Eclosedunical \nwithdescender \ eꝼ ħ\\
    ſcyꞇr he\=ı barnıno ꝼᷓ ꝼoſꞇr\textsuperscript{ı}. þa {ſ\l}o þ{\andron ꜹ} {\andron ꜹ}ðỏꝼe\\
    \k{o}ll ꝼylgıa he\=ı er ħ haꝼ\dh e ꝩı\dh {\andron ꝛ} ꞇekıꞇ. \Eunical ꝼ ħ \\
    haꝼ\dh e h\=o lı{\dh} ꝩe\'{i}ꞇꞇ. oc ꞇoc ħ aꝼ þ\textsuperscript{ı} aꞇ ꝼøða h\=o \\
    barn. {\LARGE }u l\ae ꞇr ħ b\textsuperscript{r}. he\=ı ꝼaͬ, þa {ſ\l}o b. þ͛ v͛\th a\\
    hu͛ſ ꝥ
    þıkc͛ v͛ꞇ. \Eclosedunical ꝼ mͬ heꝼ͛ ꞇeꝁ \vbart lı\dh ſ bͬ. aꞇ \\
    ꝼøþa. oc ſcyꞇr he\=ı. þa {ſ\l} ħ an ꝩ͛a lı\dh ıꞇ.\\
    \Eclosedunical ꝼ ħ l{\andron ꜹ}nar lı\dh ıꞇ barnꝼoſꞇr\textsuperscript{ı}. enꝺa ſcyꞇr he\=ı \\
    b. aꝼ ꝼoſꞇr\textsuperscript{ı}. þa {ſ\l} ħ gıallꝺa ı\=a mıꝁ ꝼe ſ\=e b.\\
    v͛\dh a aꞇ lı\dh ıꞇ v͛e v͛ꞇ. \Eclosedunical ꝼ mͬ ꞇekr ƀ. b\textsuperscript{o}ꞇꞇ ꝼᷓ \\
    ꝼoſꞇr\textsuperscript{ı} ſa e͛ \'{a}, þa {ſ\l}o epꞇ͛ {\andron ꜹ}ðỏꝼın \k{o}ll, þ{\andron ꜹ} er \\
    ꝥgaꞇ fylg\dh o, n\=ea buar be͛ ꝥ aꞇ ƀ. v͛e \textroundcap{eg} ꝩ\ae rꞇ\\
    aꞇ ꝼoſꞇr\textsuperscript{ı}no. þa {ſ\l} allꞇ rakna ꝼ\'{e}iꞇ ꝥ e͛ m; vͬ\\
    ꝼunꝺıꞇ.
    }
\end{multicols}
在这里，我们看到许多熟悉的记号，但也有一些是没有见过的或不清楚的，还有一些字形难以分辨。下面是转写本段时的一些难点：
\begin{enumerate}[itemindent=1em, label=\textbf{\arabic*}.]
    \item {\medieval ꝼøþ͛; ꝼøða}: 这两个词的不定式都是fœða，可见当时\th 和\dh 仍然混用。

    \item {\medieval {\andron ꜹ}ð{\andron ꝛ}\=o}: 这里出现了两个合体字：{\medieval {\andron ꜹ}}和{\medieval {\andron ꝛ}}，其中{\medieval {\andron ꝛ}}比较难以分辨。读者比较后面的手稿中出现的\dh 就会发现，这里的\dh 后面好像增加了一个小圆圈，这正是和{\medieval {\andron ꝛ}}合写时留下的。另外，这个词中的au对应的实际上是\k{o}，因此{\medieval {\andron ꜹ}ð{\andron ꝛ}\=o} = ǫðrum，这只能依靠读者对古诺尔斯语的知识解决。

    \item {\medieval \=me}: 这个词在m上添加了一个横杠，因此相当于缩略了一个鼻音音节。一般来说，如果被缩略的词的首字母是m的话，一般都是maðr，结合词尾的e可以判断这个词是maðr的复数与格manni.

    \item {\medieval ꝩᷓ}: 这个词比较难判断，它是vetr的复数属格vetra，我们可以看到只有首字母以及上标表示的ra保留了下来，中间的et被省略了，但数词 + vetra gamall是古诺尔斯语表示年龄的固定表达，因此可以如此缩写。

    \item {\medieval feet}: fé和后缀it，这里的正字法对判读造成一定的困难。

    \item {\medieval ꝼoſꞇr\textsuperscript{ı}}: fostri，注意最后一个i写成了上标形式以进一步节省空间。

    \item {\medieval {ſ\l}o}: skulu，注意它和{ſ\l}的区别，复数词尾由o提示。

    \item {\medieval {\andron ꜹ}ðỏꝼe}: 字典形一般记作auð-hæfi，这个词发生了进一步的音变，因此判读困难。它的字面义是“适合某人的财富”，因此表示财产。

    \item {\medieval ꝩı\dh {\andron ꝛ}}: vi\dh r，这是介词vi\dh 的另一种形式。

    \item {\medieval bͬ.}: 我们发现这里出现了一个新的缩写，代指上文重复出现的barn.

    \item {\medieval b.}: 表示búi的复数búar，这个词是“邻居”的意思，在冰岛自由邦时期，邻居常作为旁证人出庭，因此在法律文献中经常出现这个词。注意这个词和barn的缩写有所区分。

    \item {\medieval b\textsuperscript{o}ꞇꞇ}: 表示brott，现写作braut.

    \item {\medieval \textroundcap{eg}}: 表示eigi，这里的圆弧形记号只能提示读者这个词被缩略了，但并没有指出缩略的方式，读者必须结合语境和语言知识判断。

\end{enumerate}
\medskip

有一些字体也稍显独特，不过总体来说不影响阅读，例如大写的E写作\Eunical 或\Eclosedunical ，大写的N写作{\Large{\nwithdescender}} 等等。有些写法是受其他字体影响所致，例如大写字母就和安色尔体（Uncial，一种全大写字母的字体，过去曾用于抄写拉丁语和希腊语）相似，也有一些是早期哥特体特有的装饰性笔画。对于各种字体的演变历史和特点的研究不在我们的讨论范围内，因此这里不详细展开。

将正字法进一步规范化后的结果如下所示，我们根据法律文书的特点进行了分段：
\begin{paracol}{2}
    \begin{center}
        \textsc{Ómaga Bálkr}

        K 141. Um barn fóstr
    \end{center}

    Ef maðr fœðir barn ǫðrum manni ok tekr til fóstrs, ok skal hann þat barn fœða þar til er þat er XVI vetra gamalt.
    \switchcolumn
    \begin{center}
        \textsc{Dependents Section}

        Chapter 141. On child fostering
    \end{center}

    If a man raises a child for another person and undertakes fostering (for him), then he is to raise that child until he is sixteen years old.
\end{paracol}
\begin{paracol}{2}
    En ef barni verða lestir áðr þat er XVI vetra gamalt, þá skal féit aptr ganga leigulaust ok telja jamnt til hvers vetrar innstæðann.
    \switchcolumn
    But if defects befall the child before he is sixteen years old, the money (paid for fostering) is to go back without interest and it should be calculated in equal amounts for each year.
\end{paracol}

\begin{paracol}{2}
    En ef hann skýtr heim barninu frá fóstri, þá skulu þau auðhæfi ǫll fylgja heim er hann hafði viðr tekit.
    \switchcolumn
    And if he sends the child home from fostering, then all the possessions that the child owns are to go home with him.
\end{paracol}
\begin{paracol}{2}
    Ef hann hafði honum lið veitt, ok tok hann af því at fæða honum barn. Nu lætr hann barn heim fara, þá skulu búar þat virða, hvers þat þykkir vert.
    \switchcolumn
    If he has come to his aid and therefore raises the child for him, now that he has sent the child home, then the neighbors shall assess what they think the service was worth.
\end{paracol}

\begin{paracol}{2}
    Ef maðr hefir tekit til liðs barn at fœða, ok skýtr heim, þá skal hann án vera liðit.
    \switchcolumn
    If he has undertaken (the duty) of fostering and (now) sends the child home, then he is freed from the duty.
\end{paracol}

\begin{paracol}{2}
    Ef hann launar liðit barnfóstri, enda skýtr heim barn af fóstri, þá skal hann gjalda jamn-mikit fé sem búar virða at liðit væri vert.
    \switchcolumn
    If he (wishes) to reward the fostering-father for his service since he has sent the child home from fostering, then he is to pay the same amount of money that the neighbors think the service was worth.
\end{paracol}

\begin{paracol}{2}
    Ef maðr tekr barn braut frá fóstri sá er á, þá skulu eptir auðhæfin ǫll þau er þangat fylgðu, nema búar beri þat, at barn væri eigi vært at fóstrinu, þá skal allt rakna féit þat er með var fundit.
    \switchcolumn
    If the parent takes the child away from fostering, then all the resources that went there with the child are to remain unless neighbors give a verdict that the child could not live there. In that case, all the investments found for the child's keep are to be paid over.
\end{paracol}
\begin{translation*}{}
    \begin{center}
        \textbf{赡养法}

        第一百四十一章 \quad 收养子女
    \end{center}
    \AddEnumerateCounter{\chinese}{\chinese}{}
    \begin{enumerate}[label={\chinese*}、]
        \item 收养人替他人抚养子女的，应将子女抚养至16岁。

        \item 养子女不满16岁而出现残疾的，收养人应将送养人支付的财物退回。退回的钱按年等量累加，且不计利息。

        \item 收养人将养子女送回生父母时，养子女拥有的财产应一并送回。

        \item 收养人帮助他人抚养子女的，当收养人将养子女送回时，收养人的邻居应评估其抚育的成本。

        \item 收养人为帮助他人而替其抚养子女的，当其将养子女送回时，收养关系即解除。

        \item 收养人将养子女送回生父母后，若生父母愿为其支付报酬的，则应按收养人邻居所评估的价值支付酬金。

        \item 生父母将子女从养父母处带走的，则随子女送到养父母处的财物都留在养父母处，但邻居裁定养父母未能善待养子女的情况除外。此时，这些财物应退还给生父母，但生父母仍应补偿收养期间支出的费用。
    \end{enumerate}
\end{translation*}
\begin{grammar*}{}
    这篇法律文献有一定难度。值得注意的是，这段文字的主语不是很清楚，收养人、送养人以及孩子都可以用第三人称代词hann指代，这对句意的判读造成了一定困扰。

    下面分条目解释一些难点：

    \textbf{第一条}
    \begin{enumerate}
        \item fœða barn ǫðrum manni

              复数与格ǫðrum manni表示受惠的对象，与格起到的作用类似于英文的`for'.

        \item taka til

              taka能和许多介词构成丰富的含义。taka til和英语`take to'含义也有相似之处，主要表示“开始做，承担”等。
    \end{enumerate}

    \textbf{第二条}
    \begin{enumerate}
        \item lestir verða barni

              lestir是i-词干阳性名词l\k{o}ster的复数，表示“缺陷、残疾”等含义。\\verða接与格时表示“发生”，与格提示发生的对象，相当于`happen to, befall'.

        \item leigulaust

              leigu-lauss的副词形式。leiga指的是记利息的钱，它的反义词是instæða，这个词在下一句也出现了。

    \end{enumerate}

    \textbf{第四条}
    \begin{enumerate}
        \item hvers þat þykkir vert

              vert是形容词verðr `worth'的中性强变化形式，这个形容词的补语要用属格，因此疑问词采用hvers. 此外这句话省略了系动词vera，其最完整的形式是
              `hvers þat þykkir búm vera vert'.
    \end{enumerate}
    \textbf{第五条}
    \begin{enumerate}
        \item án vera

              直译为`be without'，它有时可以延申为“从...之中脱身”，类似于`free from'. 在本句中，án vera liðit的liðit就指的是抚养孩子这件事，因此这个短语可以理解为尽到了抚养的义务/责任。结合上文看，指的就是抚养关系的结束。
    \end{enumerate}
    \textbf{第六条}
    \begin{enumerate}
        \item enda

              enda是一个语义比较微妙的介词，它一般用在这样的句型中：
              \begin{quote}
                  A, enda B  其中A和B是两个分句
              \end{quote}
              根据B中动词的语气不同，enda主要表示两种含义：
              \begin{enumerate}
                  \item  虚拟式

                        表示“在A句的基础上，假定发生B句，则...”。这种情况下一般还会再接一个分句`þá ...', 此时enda表示条件，类似于`in case that, supposing that'.

                  \item 直陈式

                        此时的enda和`and'比较类似，用于承接句子。但是enda略有一点表示原因的意思，表示A、B两句在逻辑上是顺畅的。
              \end{enumerate}
              由于skýtr是直陈式，这里的enda不能理解为表条件，而应认为是大致表示承接和因果。这里的意思是，“养父母将子女送回”和“生父母支付报偿”是逻辑上关联的。另外，注意enda前后两句的主语不同，第二个分句的主语是前一个句子中的`barnfóstri'.

        \item fé sem búar virða at liðit væri vert

              注意sem引导的定语从句中的宾语从句：virða at `assess that ...'，liðit是宾语从句的主语。
    \end{enumerate}
    \textbf{第七条}
    \begin{enumerate}
        \item sá er á

              这里的á不是介词而是eiga的第三人称单数现在时，这里eiga表示的是亲缘关系：eiga barn. sá与句子的主语maðr相对应，因此maðr sá er á表示`the man who owns (the child) = the father of the child'.
    \end{enumerate}
\end{grammar*}